\documentclass[11pt]{scrartcl}


\usepackage[sexy]{evan}
\usepackage{import}
\usepackage[table]{xcolor}
\usepackage{xfp}

\newcommand{\gad}{\textcolor{yellow}{$\bigstar$}}
\newcommand{\bad}{\textcolor{red}{$\bigstar$}}

\definecolor{dcol0}{HTML}{C8E6C9}
\definecolor{dcol1}{HTML}{D4E9B3}
\definecolor{dcol2}{HTML}{E5ED9A}
\definecolor{dcol3}{HTML}{FFF59D}
\definecolor{dcol4}{HTML}{FFE082}
\definecolor{dcol5}{HTML}{FFCC80}
\definecolor{dcol6}{HTML}{FFAB91}
\definecolor{dcol7}{HTML}{F49890}
\definecolor{dcol8}{HTML}{E57373}
\definecolor{dcol9}{HTML}{D32F2F}

\makeatletter
\newcommand{\getcolorname}[1]{dcol#1}
\makeatother

\newcommand{\dif}[1]{%
    \edef\colorindex{\number\fpeval{floor(#1)}}%
    \edef\fulltext{#1}%
    \colorbox{\getcolorname{\colorindex}}{%
        \ifnum\colorindex>8
            \textbf{\textcolor{white}{\,\fulltext\,}}%
        \else
            \textbf{\textcolor{black}{\,\fulltext\,}}%
        \fi
    }%
}
% Variable para dificultad (inicial 0)
\newcommand{\thmdifficulty}{0}

% Comando para asignar dificultad antes del problema
\newcommand{\problemdiff}[1]{\renewcommand{\thmdifficulty}{#1}}

% Estilo del problema que incluye dificultad antes del título
\declaretheoremstyle[
    headfont=\color{blue!40!black}\normalfont\bfseries,
    headformat={%
      \dif{\thmdifficulty}\quad \NAME~\NUMBER\ifx\relax\EMPTY\relax\else\ \NOTE\fi
    },
    postheadspace=1em,
    spaceabove=8pt,
    spacebelow=8pt,
    bodyfont=\normalfont
]{problemstyle}

    \declaretheorem[style=problemstyle,name=Problema,sibling=theorem]{problema}
    \declaretheorem[style=problemstyle,name=Problema,numbered=no]{problema*}

\definecolor{yellow4}{RGB}{255, 206, 0}

\title {Sucesiones}

\author{Emmanuel Buenrostro}


\begin{document}
\maketitle

\section{Teoria?}


\subsection{Posibles atributos}
\begin{itemize}
\item Una sucesi\'on $\{a_n\}$ es \textit{aritmetica} si $a_n-a_{n-1}$ es constante para toda $n$.
\item Una sucesi\'on $\{a_n\}$ es \textit{geometrica} si $\frac{a_n}{a_{n-1}}$ es constante para toda $n$.
\item Una sucesi\'on $\{a_n\}$ es \textit{periodica} si $a_{n+m}=a_n$ para toda $n>C$ para algunos $m,C$.
\item Una sucesi\'on $\{a_n\}$ es \textit{acotada superiormente} si existe una constante $M$ tal que $a_n < M$ para toda $n$.
\item Una sucesi\'on $\{a_n\}$ es \textit{acotada inferiormente} si existe una constante $M$ tal que $a_n > M$ para toda $n$.
\item Una sucesi\'on $\{a_n\}$ es \textit{acotada} si esta acotada superiormente e inferiormente.
\item Una sucesi\'on $\{a_n\}$ es \textit{creciente} si $a_n>a_{n-1}$ para toda $n$.
\item Una sucesi\'on $\{a_n\}$ es \textit{decreciente} si $a_n < a_{n-1}$ para toda $n$.
\item Una sucesi\'on $\{a_n\}$ es \textit{no creciente} si $a_n \leq a_{n-1}$ para toda $n$. (Las sucesiones decrecientes son sucesiones no crecientes)
\item Una sucesi\'on $\{a_n\}$ es \textit{no decreciente} si $a_n \geq a_{n-1}$ para toda $n$. (Las sucesiones crecientes son sucesiones no decrecientes)
\item Una sucesi\'on $\{a_n\}$ es \textit{monotona} si es no creciente o no decreciente.
\end{itemize}


\subsection{¿C\'omo tienes problemas con esto?}

Bueno, usualmente los problemas te dan alguna condici\'on que cumple la sucesi\'on, puede ser
que te definan $a_n$ en terminos de los terminos anteriores, una secuencia cualquiera que cumple una condici\'on,
encontrar las sucesiones que cumplen esta condici\'on, etc. \\

Estos problemas siento que son de dos tipos, hacer cuentitas o algunas idea tipo "combinatorias", refiriendome a que tienes que 
jugar con las sucesiones, hayar alguna propiedad (por ejemplo alguna de las mencionadas anteriormente), mientras que la de hacer cuentitas 
refiriendome a cosas como encontrar una formula, calcular un valor, una recursi\'on, etc. 


\subsection{Tips}
\begin{itemize}
    \item Buscar Forma Cerrada \begin{itemize}
        \item Inducción
    \end{itemize}
    \item Desigualdades
    \item Muevele :p 
    \item Buscar algun atributo.
\end{itemize}

\newpage 
\section{Problemas}
\epigraph{Cada vez que te pase algo malo \\
"Ya te tocaba"}{}

\subsection{Nivel I}

\problemdiff{2}
\begin{problem}
%[ARML Local 2021]
Una sucesión de números $a_1, a_2, \ldots$ de números reales cumple que: 
\[a_n = na_{n-1}+(n-1)(n!(n-1)!-1)\]
para los enteros $n \geq 2$. Si $a_{2021}=(2021!+1)^2+2020!$. ¿Cuál es el valor de $a_1$?
\end{problem}

\problemdiff{3}
\begin{problem}
%[Centro 2015/2]

    Una sucesión $(a_n)$ de números reales está definida por $a_0=1$, $a_1=2015$ y para todo $n\geq1$, tenemos:
\[a_{n+1}=\frac{n-1}{n+1}a_n-\frac{n-2}{n^2+n}a_{n-1}. \]
Calcule el valor de
 $$\frac{a_1}{a_2}-\frac{a_2}{a_3}+\frac{a_3}{a_4}-\frac{a_4}{a_5}+\ldots+\frac{a_{2013}}{a_{2014}}-\frac{a_{2014}}{a_{2015}}.$$
\end{problem}

\problemdiff{3}
\begin{problem}
 %   [MEX TST 2024]
    Una secuencia de $11$ números no ceros $a_1,a_2,\ldots, a_{10},a_{11}$, cumple que $a_i=a_{i-1}a_{i+1}-1$ para $i=2,\ldots, 10$. Muestra que $a_1=a_{11}$.
\end{problem}


\problemdiff{3.5}
\begin{problem}
 %   [Centro 2021/5]
    
    Sea $n \geq 3$ un número entero y $a_1,a_2,...,a_n$ números reales positivos tales que $m$ es el menor y $M$ el mayor de estos números. Se sabe que para cualesquiera enteros distintos $1 \leq i,j,k \leq n$, si $a_i \leq a_j \leq a_k$ entonces $a_ia_k \leq a_j^2$. Demostrar que
\[ a_1a_2 \cdots a_n \geq m^2M^{n-2} \]
y determinar cuando la igualdad se mantiene.

\end{problem}


\subsection{Nivel II}

\problemdiff{4}
\begin{problem}
%    [EGMO 2023/1]
    
    Se tienen $n\geq 3$ números reales positivos $a_1,a_2,\dots,a_n$. Para cada $1\leq i\leq n$ se define $b_i=\frac{a_{i-1}+a_{i+1}}{a_i}$, donde $a_0=a_n$ y $a_{n+1}=a_1$. Suponga que para cada $1\leq i\leq n$ cada $i\leq j\leq n$ se tiene que $a_i\leq a_j$ si y sólo si $b_i\leq b_j$.  \newline 

Demuestre que $a_1=a_2=\dots =a_n$.
\end{problem}


\problemdiff{4}
\begin{problem}
 %   [EGMO 2022/4]
    
    Para cada entero positivo $n \ge 2$, determine el mayor entero positivo $N$ con la propiedad de que existen $N +1$ números reales $a_0,a_1,\dots,a_N$ tales que
 \begin{itemize} 
 \item  $a_0+a_1=-\frac 1n$, y
 \item  $(a_k+a_{k-1})(a_k+a_{k+1})=a_{k-1}-a_{k+1}$ para todo $1\leq k \leq N-1$.
 \end{itemize} 
    
\end{problem}

\problemdiff{4}
\begin{problem}
 %   [MEX EGMO TST 2025]
    ¿Existen 2024 números reales distintos de cero $a_1,a_2, \ldots, a_{2024}$ tales que:
    \[\sum_{i=1}^{2024}\left(a_i^2+\frac{1}{a_{i}^2}\right) + 2\sum_{i=1}^{2024}\frac{a_i}{a_{i+1}}+2024=2\sum_{i=1}^{2024}\left(a_i+\frac{1}{a_i}\right)?\]
\end{problem}

\problemdiff{4.5}
\begin{problem}
 %   [Canadian Training Camp]
    Las sucesiones $a_n$ y $b_n$ cumplen que para cada entero positivo $n$ 
    \[ a_n > 0,\qquad\ b_n>0,\qquad\ a_{n+1}=a_n+\dfrac{1}{b_n},\qquad\ b_{n+1} = b_n+\dfrac{1}{a_n}. \]
    Prueba que $a_{50}+b_{50}>20$.
\end{problem}

\subsection{Nivel III}

\problemdiff{5}
\begin{problem}
 %   [EGMO 2024/4]
    Para una sucesión $a_1 < a_2 < \ldots < a_n$ de enteros, un par $(a_i, a_j)$ con $1 \leq i < j \leq n$ es llamada \textit{interesante} si existe un par $(a_k, a_l)$ de enteros con $1\leq k < l \leq n$ tal que:
    $$\frac{a_l-a_k}{a_j-a_i}=2.$$
    Para cada $n \geq 3$, encuentra el número de pares \textit{interesantes} en una sucesión de tamaño $n$.
\end{problem}

\problemdiff{5}
\begin{problem}
 %   [OMM 2011/3]
    
    Sea $n$ un entero positivo. Encuentra todas las soluciones $(a_1, a_2,\dots , a_n)$ de números reales
que satisfacen el siguiente sistema de $n$ ecuaciones:

\[a_1^2 + a_1 - 1 = a_2\]
\[ a_2^2 + a_2 - 1 = a_3\]
\[\vdots \]
\[a_{n}^2 + a_n - 1 = a_1.\]
 
\end{problem}


\problemdiff{5}
\begin{problem}
 %   [OMM 2020/6]
    
    Sea $n\ge 2$ un número entero positivo. Sean $x_1, x_2, \dots, x_n$ números reales no nulos que satisfacen la ecuación
\[\left(x_1+\frac{1}{x_2}\right)\left(x_2+\frac{1}{x_3}\right)\dots\left(x_n+\frac{1}{x_1}\right)=\left(x_1^2+\frac{1}{x_2^2}\right)\left(x_2^2+\frac{1}{x_3^2}\right)\dots\left(x_n^2+\frac{1}{x_1^2}\right). \]
Encuentra todos los valores posibles de $x_1, x_2, \dots, x_n$.
\end{problem}

\problemdiff{5}
\begin{problem}
 %   [IMOSL 2015/A1]
    Supón que una sucesión $a_1,a_2,\ldots$ de números reales positivos satisface\[a_{k+1}\geq\frac{ka_k}{a_k^2+(k-1)}\]para cada entero positivo $k$. Demuestra que $a_1+a_2+\ldots+a_n\geq n$ para cada $n\geq2$.
\end{problem}


\problemdiff{5.5}
\begin{problem}
  %  [EGMO 2025/2]
    Una sucesión infinita creciente  $a_1 < a_2 < a_3 < \ldots$ de enteros positivos es llamada \textit{central} si para cada entero positivo $n$, la media aritmetica de los primeros $a_n$ terminos es igual a $a_n$. \\
    Prueba que existe una sucesión infinita $b_1, b_2, b_3,\ldots$ de enteros positivos tal que para cada sucesión central $a_1, a_2, a_3, \ldots$ hay infinitos enteros positivos $n$ tal que $a_n=b_n$. 
\end{problem}

\subsection{Nivel IV}

\problemdiff{6}
\begin{problem}
  %  [EGMO 2020/2]
    
    Encontrar todas las listas $(x_1, x_2, \ldots, x_{2020})$ de números reales no negativos tales que se satisfagan las tres condiciones siguientes:

    $x_1 \le x_2 \le \ldots \le x_{2020}$; $\qquad$  
    $x_{2020} \le x_1 + 1$;  $\qquad$ 
    Existe una permutación $(y_1, y_2, \ldots, y_{2020})$ de $(x_1, x_2, \ldots, x_{2020})$ tal que

\[\sum_{I = 1}^{2020} ((x_i + 1)(y_i + 1))^2 = 8 \sum_{i = 1}^{2020} x_i^3\]


Una permutación de una lista es una lista de la misma longitud, con los mismos elementos pero en un
orden cualquiera. Por ejemplo, $(2, 1, 2)$ es una permutación de $(1, 2, 2)$, y ambas son permutaciones
de $(2, 2, 1)$. En particular cualquier lista es una permutación de ella misma.
    
\end{problem}


\problemdiff{6}
\begin{problem}
  %  [IMO 2018/2]
    Encuentra todos los enteros $n \geq 3$ para los cuales existen números reales $a_1, a_2, \dots a_{n + 2}$ que satisfacen $a_{n + 1} = a_1$, $a_{n + 2} = a_2$ y
$$a_ia_{i + 1} + 1 = a_{i + 2},$$for $i = 1, 2, \dots, n$.
\end{problem}

\problemdiff{7}
\begin{problem}
 %   [IMOSL 2024/A5]
    Encuentra todas las sucesiones periodicas $a_1,a_2, \ldots$ de reales tales que lo siguiente se cumple para todo $n \geq 1$.
    $$a_{n+2}+a_{n}^2=a_n+a_{n+1}^2\quad\text{y}\quad |a_{n+1}-a_n|\leqslant 1.$$
\end{problem}

\newpage
\section{Problemas Parte II}
\epigraph{Vemos jaja}

\subsection{Nivel I}

\begin{problem}
 %   [ASHME 1999/20]
    La sucesión $a_1, a_2, \ldots, a_n$ cumple que $a_1=19, a_9=99$ y para todo $n\geq 3$, $a_n$ es el promedio de los primeros $n-1$ terminos. Calcula $a_2$. 
\end{problem}


\begin{problem}
  %  [Baltic Way 2025/1]
    Does there exist a sequence of integers $a_1, a_2,\dots$ such that for each integer $d\neq 0$ there are exactly $2025$ distinct pairs of indices $(i,j)$ for which $a_i - a_j = d$?
\end{problem}

\subsection{Nivel II}
\begin{problem}
  %  [EGMO 2016/1]
    Sean $n$ un entero positivo impar, y $x_1, \dots, x_n$ números reales no negativos. Muestra
que
\[ \min_{i=1,\ldots,n} (x_i^2+x_{i+1}^2) \leq \max_{j=1,\ldots,n} (2x_jx_{j+1}) \]
donde $x_{n+1}= x_1$.
    
\end{problem}


\begin{problem}
   % [Ibero 2009/5]
    
    La sucesión $a_n$ está definida por $a_1 = 1$, $a_{2k} = 1 + a_k$ y $a_{2k+1} = \frac{1}{a_{2k}}$, para todo entero $k \geq 1$. \newline 
Demostrar que todo número racional positivo aparece exactamente una vez en esta sucesión.

    
\end{problem}


\begin{problem}
    %[IMOSL 2013/A1]
    Sea $n$ un entero positivo y sean $a_1, \ldots, a_{n-1} $ números reales arbitrarios. Definimos las sucesiones $u_0, \ldots, u_n $ y $v_0, \ldots, v_n $ inductivamente por $u_0 = u_1  = v_0 = v_1 = 1$, y $u_{k+1} = u_k + a_k u_{k-1}$, $v_{k+1} = v_k + a_{n-k} v_{k-1}$ para $k=1, \ldots, n-1.$

Demuestra que $u_n = v_n.$
\end{problem}


\subsection{Nivel III}

\begin{problem}

  %  [IMO 2014/1]
    
    Sea $a_0 < a_1 < a_2 < \dots$ una sucesión infinita de números enteros positivos. Demostrar que existe un único entero $n \geq 1$ tal que
\[a_n < \frac{a_0+a_1+a_2+\cdots+a_n}{n} \leq a_{n+1}.\]
    
\end{problem}


\begin{problem}
 %   [IMO SL 2006/A1]
     Una sucesión de números reales $ a_{0},\ a_{1},\ a_{2},\dots$ se define por la fórmula
\[ a_{i + 1} = \left\lfloor a_{i}\right\rfloor\cdot \left\langle a_{i}\right\rangle\qquad\text{para}\quad i\geq 0;
\]aquí $a_0$ es un número real arbitrario, $\lfloor a_i\rfloor$ denota el mayor entero que no excede $a_i$, y $\left\langle a_i\right\rangle=a_i-\lfloor a_i\rfloor$. Muestra que $a_i=a_{i+2}$ para todo $i$ suficientemente grande.
\end{problem}

\begin{problem}
%    [IMOSL 2014/A2]
     Se define la función $f:(0,1)\to (0,1)$ por\[\displaystyle f(x) = \left\{ \begin{array}{lr} x+\frac 12 & \text{si}\ \  x < \frac 12\\ x^2 & \text{si}\ \  x \ge \frac 12 \end{array} \right.\]Sean $a$ y $b$ dos números reales tales que $0 < a < b < 1$. Se definen las secuencias $a_n$ y $b_n$ por $a_0 = a, b_0 = b$, y $a_n = f( a_{n -1})$, $b_n = f (b_{n -1} )$ para $n > 0$. Muestra que existe un entero positivo $n$ tal que \[(a_n - a_{n-1})(b_n-b_{n-1})<0.\]
\end{problem}


\begin{problem}
 %   [IMO 2007/1]
    Sean $ a_{1}$, $ a_{2}$, $ \ldots$, $ a_{n}$ números reales. Para cada $ i$, $ (1 \leq i \leq n )$, se define
\[ d_{i} = \max \{ a_{j}\mid 1 \leq j \leq i \} - \min \{ a_{j}\mid i \leq j \leq n \}
\]
y sea $ d = \max \{d_{i}\mid 1 \leq i \leq n \}$.

(a) Muestra que, para cualesquiera números reales $ x_{1}\leq x_{2}\leq \cdots \leq x_{n}$,
\[ \max \{ |x_{i} - a_{i}| \mid 1 \leq i \leq n \}\geq \frac {d}{2}. \quad \quad (*)
\]
(b) Muestra que existen números reales $ x_{1}\leq x_{2}\leq \cdots \leq x_{n}$ tales que se da la igualdad en (*).
\end{problem}


\begin{problem}
 %   [IMOSL 2022/A1]
    Let $(a_n)_{n\geq 1}$ be a sequence of positive real numbers with the property that
$$(a_{n+1})^2 + a_na_{n+2} \leq a_n + a_{n+2}$$for all positive integers $n$. Show that $a_{2022}\leq 1$.
\end{problem}
\begin{problem}
  %  [EGMO 2015/4]
    
    Determina si existe una sucesión infinita $a_1, a_2, a_3, \dots$ de enteros positivos que satisface
la igualdad

\[a_{n+2}=a_{n+1}+\sqrt{a_{n+1}+a_{n}} \]

para todo entero positivo $n$.
    
\end{problem}

\begin{problem}
  %  [Baltic Way 2025/2]
    Let $a_1, a_2, \dots$ be a sequence of real numbers such that
$$ \{a_n\}\lfloor {a_{n+1}} \rfloor = \lfloor{a_n} \rfloor \{a_{n+1}\} $$for each positive integer $n$.
Prove that there is a real number $\lambda$ such that $\{a_m\}\lfloor{a_m} \rfloor = \lambda$ for infinitely many positive integers $m$.
Remark: For a real number \(x\), \(\lfloor x \rfloor\) denotes the greatest integer that does not exceed \(x\), and ${\{x\}=x-\lfloor x \rfloor}$.
\end{problem}

\begin{problem}
%    [EGMO 2022/3]
    
    Se dice que una sucesión infinita de enteros positivos $a_1, a_2, \dots$ es húngara si
 \begin{itemize} 
 \item  $a_1$ es un cuadrado perfecto, y
 \item  para todo entero $n\ge 2$, $a_n$ es el menor entero positivo tal que \[na_1+(n-1)a_2+\dots+2a_{n-1}+a_n\]
es un cuadrado perfecto.
 \end{itemize} 
Pruebe que si $a_1, a_2, \dots$ es una sucesión húngara, entonces existe un entero positivo $k$ tal que $a_n=a_k$ para todo entero $n\geq k$.
    
\end{problem}


\subsection{Nivel IV}

\begin{problem}
 %   [IMO SL 2013/A4]
    Sea $n$ un entero positivo, y considera la sucesión $a_1 , a_2 , \cdots , a_n $ de enteros positivos. Esta se extiende de forma periódica a $a_1 , a_2 , \cdots $ definiendo $a_{n+i} = a_i $ para cada $i \ge 1$. Si\[a_1 \le a_2 \le \cdots \le a_n \le a_1 +n  \]y \[a_{a_i } \le n+i-1 \quad\text{for}\quad i=1,2,\cdots, n, \]muestra que\[a_1 + \cdots +a_n \le n^2. \]
\end{problem}

\begin{problem}
 %   [IMO SL 2018/A4]
    Sea $a_0,a_1,a_2,\dots $ una sucesión de números reales tales que $a_0=0, a_1=1,$ y para cada $n\geq 2$ existe algún $1 \leq k \leq n$ que satisface\[ a_n=\frac{a_{n-1}+\dots + a_{n-k}}{k}. \]Encuentra el mayor valor posible de $a_{2018}-a_{2017}$.
\end{problem}

\begin{problem}
   % [IMO SL 2023/A5]
    Let $a_1,a_2,\dots,a_{2023}$ be positive integers such that
$a_1,a_2,\dots,a_{2023}$ is a permutation of $1,2,\dots,2023$, and
$|a_1-a_2|,|a_2-a_3|,\dots,|a_{2022}-a_{2023}|$ is a permutation of $1,2,\dots,2022$.
Prove that $\max(a_1,a_{2023})\ge 507$.
\end{problem}

\begin{problem}
 %   [IMO 2023/3]
    For each integer $k\geq 2$, determine all infinite sequences of positive integers $a_1$, $a_2$, $\ldots$ for which there exists a polynomial $P$ of the form\[ P(x)=x^k+c_{k-1}x^{k-1}+\dots + c_1 x+c_0, \]where $c_0$, $c_1$, \dots, $c_{k-1}$ are non-negative integers, such that\[ P(a_n)=a_{n+1}a_{n+2}\cdots a_{n+k} \]for every integer $n\geq 1$
\end{problem}



\end{document}